\documentclass[12pt,a4paper]{article}

\usepackage[utf8]{inputenc}
\usepackage[T1]{fontenc}
\usepackage[english]{babel}
\usepackage[a4paper,margin=2.2cm]{geometry}
\usepackage{xcolor}
\usepackage{graphicx}
\usepackage{float}
\usepackage{booktabs}
\usepackage{hyperref}
\usepackage{listings}
\usepackage{upquote}
\usepackage{fancyhdr}
\usepackage{enumitem}
\usepackage{array}

\definecolor{codebackground}{RGB}{245,247,250}
\definecolor{codecomment}{RGB}{80,120,80}
\definecolor{codekeyword}{RGB}{0,70,150}
\definecolor{codestring}{RGB}{150,50,20}
\definecolor{linkblue}{RGB}{0,75,140}

\hypersetup{
    colorlinks=true,
    linkcolor=linkblue,
    urlcolor=linkblue,
    pdftitle={TP7 - PHP PDO CRUD Application}
}

\lstdefinelanguage{PHP}{
    morekeywords={class,function,if,else,elseif,try,catch,throw,new,require,require_once,
        return,foreach,while,as,isset,empty,exit,echo,extends,public,private,protected,
        static,true,false,null},
    sensitive=true,
    morecomment=[l]{//},
    morecomment=[s]{/*}{*/},
    morestring=[b]"
}

\lstdefinelanguage{CSS}{
    morekeywords={color,background,background-color,display,flex,font-family,
        margin,padding,width,height,border,border-radius,box-shadow},
    morecomment=[s]{/*}{*/},
    morestring=[b]"
}

\lstdefinelanguage{JSON}{
    morestring=[b]",
    morecomment=[l]{//}
}

\lstdefinelanguage{Bash}{
    morekeywords={php,composer,install,localhost},
    morecomment=[l]{\#},
    morestring=[b]"
}

\lstset{
    language=PHP,
    backgroundcolor=\color{codebackground},
    basicstyle=\ttfamily\footnotesize,
    keywordstyle=\color{codekeyword}\bfseries,
    commentstyle=\color{codecomment}\itshape,
    stringstyle=\color{codestring},
    numbers=left,
    numberstyle=\tiny\color{gray},
    numbersep=8pt,
    frame=single,
    rulecolor=\color{gray!45},
    breaklines=true,
    breakatwhitespace=false,
    showstringspaces=false,
    tabsize=4,
    columns=fullflexible
}

\pagestyle{fancy}
\fancyhf{}
\lhead{TP7 -- PHP / MySQL}
\rhead{PDO CRUD application}
\cfoot{\thepage}

\newcommand{\screenshotplaceholder}[2]{%
    \begin{figure}[H]
        \centering
        \fbox{%
            \begin{minipage}[c][5.0cm][c]{0.88\textwidth}
                \centering
                \large\textbf{Screenshot placeholder}\\[0.4cm]
                \texttt{#1}\\[0.4cm]
                \normalsize Replace this box with the corresponding screenshot.
            \end{minipage}%
        }
        \caption{#2}
    \end{figure}
}

\newcommand{\insertScreenshot}[3]{%
    \begin{figure}[H]
        \centering
        \includegraphics[width=0.92\textwidth]{#1}
        \caption{#2}
        \label{#3}
    \end{figure}
}

\title{\textbf{TP7 -- PHP PDO CRUD Application}\\
\large User management, data validation and visualization}
\author{Student: Walid Amjahdi\\[0.2cm]S5 ISI -- PHP}
\date{\today}

\begin{document}

\begin{titlepage}
    \centering
    \vspace*{2cm}
    {\Large\textsc{S5 ISI -- PHP}}\\[1.2cm]
    {\Huge\bfseries TP7}\\[0.5cm]
    {\LARGE PHP PDO CRUD Application}\\[1.5cm]
    \rule{\textwidth}{0.5pt}\\[0.5cm]
    {\large User management with MySQL, validation and JpGraph}\\[0.5cm]
    \rule{\textwidth}{0.5pt}\\[2cm]
    \begin{tabular}{>{\bfseries}rl}
        Student: & Walid Amjahdi\\
        Course: & PHP\\
        Technologies: & PHP, PDO, MySQL, Composer, JpGraph\\
        Academic year: & 2026--2027
    \end{tabular}
    \vfill
    {\large \today}
\end{titlepage}

\tableofcontents
\newpage

\section{Introduction}

This practical work implements a small user-management web application in PHP.
The application uses PDO to connect to a MySQL database and provides the four
fundamental CRUD operations:

\begin{itemize}[nosep]
    \item create a user;
    \item display the list of users;
    \item modify an existing user;
    \item delete a user.
\end{itemize}

The project also includes a command-line data seeder based on Faker and a
JpGraph chart representing users' birth years according to their CIN values.

\section{Objectives}

The objectives of TP7 are:

\begin{enumerate}[nosep]
    \item configure a reusable PDO database connection;
    \item manipulate MySQL data from PHP;
    \item use prepared statements for write operations;
    \item validate form data before insertion or update;
    \item present database records in an HTML table;
    \item add visual feedback after each CRUD operation;
    \item generate test data and visualize it in a chart.
\end{enumerate}

\section{Project structure}

\begin{table}[H]
\centering
\begin{tabular}{>{\ttfamily}p{0.25\textwidth}p{0.64\textwidth}}
\toprule
\textnormal{File} & \textnormal{Responsibility}\\
\midrule
config.php & Creates the PDO connection and configures error handling.\\
index.php & Reads and displays all users; embeds the chart.\\
add.php & Displays the creation form and inserts a validated user.\\
modify.php & Loads one user, displays its form and updates it.\\
delete.php & Deletes a user identified by its ID.\\
seed.php & Generates fake users with Faker from the command line.\\
chart.php & Reads birth dates and produces a JpGraph line chart.\\
index.css & Defines the dark user interface and form/table styles.\\
composer.json & Declares Faker and JpGraph dependencies.\\
\bottomrule
\end{tabular}
\caption{Main project files}
\end{table}

\section{Database design}

The PHP code expects a MySQL database named \texttt{php\_crud} containing a
\texttt{users} table. A compatible schema is shown below:

\begin{lstlisting}[language=SQL,caption={SQL schema for the users table}]
CREATE DATABASE IF NOT EXISTS php_crud
    CHARACTER SET utf8mb4
    COLLATE utf8mb4_unicode_ci;

USE php_crud;

CREATE TABLE IF NOT EXISTS users (
    id INT UNSIGNED AUTO_INCREMENT PRIMARY KEY,
    cin VARCHAR(8) NOT NULL UNIQUE,
    full_name VARCHAR(150) NOT NULL,
    tel VARCHAR(10) NOT NULL,
    email VARCHAR(255) NOT NULL UNIQUE,
    birth DATE NOT NULL
);
\end{lstlisting}

\textbf{Note:} the database account password is intentionally not reproduced
in this report. It must be supplied locally in \texttt{config.php} or through
a safer environment-based configuration.

\section{Application implementation}

\subsection{Database connection}

The connection uses the PDO MySQL driver, UTF-8 encoding, exception-based
errors, associative fetches and native prepared statements. This centralizes
the database configuration and allows the other pages to reuse the
\texttt{\$pdo} object.

\begin{lstlisting}[caption={Documented version of config.php}]
<?php
$dsn = "mysql:host=localhost;dbname=php_crud;charset=utf8mb4";

try {
    // Replace DB_PASSWORD with the local password; do not commit credentials.
    $pdo = new PDO($dsn, "root", "DB_PASSWORD", [
        PDO::ATTR_ERRMODE => PDO::ERRMODE_EXCEPTION,
        PDO::ATTR_DEFAULT_FETCH_MODE => PDO::FETCH_ASSOC,
        PDO::ATTR_EMULATE_PREPARES => false,
    ]);
} catch (PDOException $e) {
    error_log("Database connection failed: " . $e->getMessage());
    throw new RuntimeException("Unable to connect to the database.", 0, $e);
}
\end{lstlisting}

\subsection{Listing users}

\texttt{index.php} executes a \texttt{SELECT} query, fetches all records, and
renders them in an HTML table. Values are escaped with
\texttt{htmlspecialchars()} before being written to the page. The page also
displays success or error messages and links to the creation, modification and
deletion actions.

\begin{lstlisting}[caption={Core listing query}]
require_once "config.php";
$stmt = $pdo->query(
    "SELECT id, cin, full_name, tel, email, birth
     FROM users ORDER BY id DESC"
);
$users = $stmt->fetchAll();
\end{lstlisting}

\subsection{Creating a user}

The creation form collects the CIN, full name, telephone number, email and
birth date. The submitted values are trimmed, checked with regular expressions,
and inserted using a prepared statement.

\begin{lstlisting}[caption={Core insertion logic}]
$stmt = $pdo->prepare(
    "INSERT INTO users (cin, full_name, tel, email, birth)
     VALUES (?, ?, ?, ?, ?)"
);
$stmt->execute([$cin, $full_name, $tel, $email, $birth]);
header("Location: index.php?success=1");
exit();
\end{lstlisting}

\subsection{Updating a user}

The modification page first loads the selected record using the ID from the
query string. After the same validation process as the creation page, it
executes an \texttt{UPDATE} prepared statement and redirects to the list.

\begin{lstlisting}[caption={Core update logic}]
$stmt = $pdo->prepare(
    "UPDATE users
     SET cin = ?, full_name = ?, tel = ?, email = ?, birth = ?
     WHERE id = ?"
);
$stmt->execute([$cin, $full_name, $tel, $email, $birth, $id]);
\end{lstlisting}

\subsection{Deleting a user}

The deletion page receives the user ID, executes a parameterized
\texttt{DELETE}, then redirects to the list with a confirmation flag.

\begin{lstlisting}[caption={Core deletion logic}]
$stmt = $pdo->prepare("DELETE FROM users WHERE id = ?");
$stmt->execute([$id]);
header("Location: index.php?deleted=1");
exit();
\end{lstlisting}

\subsection{Input validation}

The application checks the following formats before insertion or update:

\begin{table}[H]
\centering
\begin{tabular}{>{\ttfamily}p{0.21\textwidth}p{0.67\textwidth}}
\toprule
\textnormal{Field} & \textnormal{Rule}\\
\midrule
cin & Two uppercase letters followed by six digits, for example \texttt{AB123456}.\\
email & Basic email pattern containing an address and a domain.\\
tel & Moroccan-style number beginning with \texttt{05}, \texttt{06} or \texttt{07}.\\
birth & HTML date input; the database stores a \texttt{DATE}.\\
\bottomrule
\end{tabular}
\caption{Validation rules implemented in add.php and modify.php}
\end{table}

\section{Data generation and visualization}

\subsection{Faker seeder}

\texttt{seed.php} generates between 1 and 100 fake users. It uses the French
locale, inserts data inside a transaction, and rolls back if an unexpected
error occurs. The command is:

\begin{lstlisting}[language=Bash,caption={Installing dependencies and seeding data}]
composer install
php seed.php 15
\end{lstlisting}

\subsection{JpGraph chart}

\texttt{chart.php} retrieves each user's CIN and birth date, converts the date
to a year, and generates a line chart. The chart is embedded in
\texttt{index.php} with:

\begin{lstlisting}[language=HTML,caption={Chart inclusion in index.php}]
<section class="chart-card">
    <h2>CIN and birth year</h2>
    <p class="muted">Each point represents one registered user.</p>
    <img src="chart.php" alt="Chart showing CIN and birth year">
</section>
\end{lstlisting}

\section{User interface}

The stylesheet provides a dark theme, responsive page width, styled forms,
table rows, feedback messages and distinct action buttons. The interface
contains the following principal screens.

\screenshotplaceholder{screenshots/01-users-list.png}{Users list with CRUD actions and the chart}
\screenshotplaceholder{screenshots/02-add-user.png}{Add-user form}
\screenshotplaceholder{screenshots/03-modify-user.png}{Modify-user form populated with an existing user}
\screenshotplaceholder{screenshots/04-success-message.png}{Success message after a CRUD operation}
\screenshotplaceholder{screenshots/05-chart.png}{Generated CIN and birth-year chart}

\paragraph{Adding real screenshots.}
Save screenshots with the names shown above inside the \texttt{screenshots/}
folder. Then replace a placeholder, for example:

\begin{lstlisting}[language={[LaTeX]TeX},caption={Replacing a placeholder with an image}]
\insertScreenshot{screenshots/01-users-list.png}
    {Users list with CRUD actions and the chart}{fig:users-list}
\end{lstlisting}

\section{Execution procedure}

\begin{enumerate}
    \item Start the MySQL server.
    \item Create the \texttt{php\_crud} database and the \texttt{users} table.
    \item Set the local database password in \texttt{config.php}.
    \item From the TP7 directory, install PHP dependencies:
\begin{lstlisting}[language=Bash]
composer install
\end{lstlisting}
    \item Generate optional demonstration data:
\begin{lstlisting}[language=Bash]
php seed.php 15
\end{lstlisting}
    \item Start the PHP development server:
\begin{lstlisting}[language=Bash]
php -S localhost:8000
\end{lstlisting}
    \item Open \url{http://localhost:8000/index.php}.
\end{enumerate}

\section{Functional tests}

\begin{table}[H]
\centering
\begin{tabular}{p{0.09\textwidth}p{0.45\textwidth}p{0.28\textwidth}}
\toprule
\textbf{ID} & \textbf{Test} & \textbf{Expected result}\\
\midrule
T1 & Open the home page with a working database. & Users are displayed.\\
T2 & Submit a valid new user. & User is inserted and a success message appears.\\
T3 & Submit an invalid CIN. & The request is rejected.\\
T4 & Submit an invalid email or telephone number. & The request is rejected.\\
T5 & Modify an existing user with valid data. & The record is updated.\\
T6 & Delete an existing user. & The record disappears from the list.\\
T7 & Run \texttt{php seed.php 15}. & Fifteen fake users are inserted.\\
T8 & Open the chart after adding users. & A birth-year chart is generated.\\
\bottomrule
\end{tabular}
\caption{Recommended functional test plan}
\end{table}

\section{Security and quality observations}

\begin{itemize}
    \item Prepared statements are used for insertion, update and deletion,
    reducing SQL-injection risk.
    \item Output values are escaped before HTML rendering.
    \item PDO exceptions are logged and user-facing messages avoid exposing
    database details.
    \item The database password should be moved to an environment variable
    before deployment.
    \item The delete action is currently triggered by a link; a production
    version should use a POST request with CSRF protection and confirmation.
    \item Validation can be strengthened with server-side length checks,
    database constraints and a dedicated validation helper.
\end{itemize}

\section{Conclusion}

This TP7 project demonstrates a complete PHP/MySQL CRUD workflow using PDO.
It combines database access, server-side validation, HTML rendering, dependency
management, generated test data and graphical reporting. The structure remains
simple and readable, while the use of prepared statements, output escaping and
transactional seeding provides a solid foundation for further improvements.

\appendix
\section{Complete source code}

All relevant source code is embedded directly in this report. No external
source files are required by the LaTeX document. The database password is
redacted in the configuration listing.

\subsection{\texttt{config.php}}
\begin{lstlisting}[language=PHP]
<?php
$dsn = "mysql:host=localhost;dbname=php_crud;charset=utf8mb4";

try {
    // The local database password is intentionally redacted.
    $pdo = new PDO($dsn, "root", "DB_PASSWORD", [
        PDO::ATTR_ERRMODE => PDO::ERRMODE_EXCEPTION,
        PDO::ATTR_DEFAULT_FETCH_MODE => PDO::FETCH_ASSOC,
        PDO::ATTR_EMULATE_PREPARES => false,
    ]);
} catch (PDOException $e) {
    error_log("Database connection failed: " . $e->getMessage());
    throw new RuntimeException("Unable to connect to the database.", 0, $e);
}
\end{lstlisting}

\subsection{\texttt{index.php}}
\begin{lstlisting}[language=PHP]
<?php
$users = [];
$errorMessage = null;

try {
    require_once "config.php";
    $stmt = $pdo->query("SELECT id, cin, full_name, tel, email, birth FROM users ORDER BY id DESC");
    $users = $stmt->fetchAll();
} catch (Throwable $e) {
    error_log("Unable to load users: " . $e->getMessage());
    $errorMessage = "Users could not be loaded. Please try again later.";
}
?>

<!DOCTYPE html>
<html lang="en">
<head>
    <meta charset="UTF-8">
    <meta name="viewport" content="width=device-width, initial-scale=1.0">
    <title>Users List</title>
    <link rel="stylesheet" href="index.css"/>
</head>
<body>
    <?php if ($errorMessage !== null): ?>
        <div class="message error"><?= htmlspecialchars($errorMessage, ENT_QUOTES, 'UTF-8') ?></div>
    <?php elseif (isset($_GET['success']) && $_GET['success'] == 1): ?>
        <div class="message success">User added successfully!</div>
    <?php elseif (isset($_GET['updated']) && $_GET['updated'] == 1): ?>
        <div class="message success">User updated successfully!</div>
    <?php elseif (isset($_GET['deleted']) && $_GET['deleted'] == 1): ?>
        <div class="message success">User deleted successfully!</div>
    <?php endif; ?>

    <div class="add-div">
        <h1>Users List</h1> 
        <a href="add.php" class="add-user-link">
            Add User
        </a>
    </div>  

    <table>
        <thead>
            <tr>
                <th>ID</th>
                <th>CIN</th>
                <th>Full Name</th>
                <th>Tel</th>
                <th>Email</th>
                <th>Birthdate</th>
                <th>Actions</th>
            </tr>
        </thead>
        <tbody>
        <?php if (empty($users)): ?>
            <tr><td colspan="7" style="text-align: center;">No users found yet.</td></tr>
        <?php else: ?>
                <?php foreach ($users as $user): ?>
                <tr>
                    <td><?= htmlspecialchars((string) $user['id'], ENT_QUOTES, 'UTF-8') ?></td>
                    <td><?= htmlspecialchars($user['cin'], ENT_QUOTES, 'UTF-8') ?></td>
                    <td><?= htmlspecialchars($user['full_name'], ENT_QUOTES, 'UTF-8') ?></td>
                    <td><?= htmlspecialchars($user['tel'], ENT_QUOTES, 'UTF-8') ?></td>
                    <td><?= htmlspecialchars($user['email'], ENT_QUOTES, 'UTF-8') ?></td>
                    <td><?= htmlspecialchars($user['birth'], ENT_QUOTES, 'UTF-8') ?></td>
                    <td class="action">
                        <a href="modify.php?id=<?= (int) $user['id'] ?>" class="modify">Modify</a>
                        <a href="delete.php?id=<?= (int) $user['id'] ?>" class="delete">Delete</a>
                    </td>
                </tr>
                <?php endforeach; ?>
            <?php endif; ?>
        </tbody>
    </table>

    <section class="chart-card">
        <h2>CIN and birth year</h2>
        <p class="muted">Each point represents one registered user.</p>
        <img src="chart.php" alt="Chart showing CIN and birth year">
    </section>

</body>
</html>
\end{lstlisting}

\subsection{\texttt{add.php}}
\begin{lstlisting}[language=PHP]
<?php
$errorMessage = null;
try {
    require_once "config.php";
} catch (Throwable $e) {
    error_log("Add page database error: " . $e->getMessage());
    $errorMessage = "The database is currently unavailable.";
}
?>

<!DOCTYPE html>
<html lang="en">
<head>
    <meta charset="UTF-8">
    <meta name="viewport" content="width=device-width, initial-scale=1.0">
    <title>Add User</title>
    <link rel="stylesheet" href="index.css"/>
</head>
<body>
    <?php if ($errorMessage !== null): ?>
        <div class="message error"><?= htmlspecialchars($errorMessage, ENT_QUOTES, 'UTF-8') ?></div>
    <?php else: ?>
    <h1>Add User Infos</h1>
    <form method="POST">
        <label for="CIN">CIN :</label>
        <input type="text" name="cin" id="CIN" placeholder="AB123456" required>

        <label for="name">Full Name :</label>
        <input type="text" name="full_name" id="name" required>

        <label for="tel">Tel :</label>
        <input type="text" name="tel" id="tel" placeholder="0612345678" required>

        <label for="email">Email :</label>
        <input type="email" name="email" id="email" required>

        <label for="birth">Birthdate :</label>
        <input type="date" name="birth" id="birth" required>
        <br><br>
        <button class="add-user" type="submit" name="submit">Save User</button>
    </form>
    <?php endif; ?>
</body>
</html>

<?php

if (isset($_POST['submit']) && $_SERVER["REQUEST_METHOD"] == "POST")
{
    $formData = [];
    foreach ($_POST as $key => $value) 
    {
        $formData[$key] = trim($value);
    }
    
    $cin      = $formData["cin"] ?? '';
    $full_name = $formData["full_name"] ?? '';
    $tel      = $formData["tel"] ?? '';
    $email    = $formData["email"] ?? '';
    $birth    = $formData["birth"] ?? '';
    
    if (preg_match("/^[A-Z]{2}[0-9]{6}$/", $cin) == 0) 
    {
        exit("Invalid CIN format (e.g., AB123456)");
    }
    if (preg_match("/^[A-Za-z0-9.]+@[A-Za-z]+\.[a-z.]+$/", $email) == 0) 
    {
        exit("Invalid Email format");
    }
    if (preg_match("/^0[567]{1}[0-9]{8}$/", $tel) == 0) 
    {
        exit("Invalid Tel format (e.g., 0612345678)");
    }
    
    try {
        $stmt = $pdo->prepare("INSERT INTO users (cin, full_name, tel, email, birth) VALUES (?, ?, ?, ?, ?)");
        $stmt->execute([$cin, $full_name, $tel, $email, $birth]);
    
        header("Location: index.php?success=1");
        exit();
    
    } catch (Throwable $e) {
        error_log("Unable to add user: " . $e->getMessage());
        exit("Unable to save the user. Please try again.");
    }
}
?>
\end{lstlisting}

\subsection{\texttt{modify.php}}
\begin{lstlisting}[language=PHP]
<?php
try {
    require_once "config.php";
} catch (Throwable $e) {
    error_log("Modify page database error: " . $e->getMessage());
    exit("The database is currently unavailable.");
}

if (!isset($_GET['id'])) {
    header("Location: index.php");
    exit();
}

$id = $_GET['id'];

$stmt = $pdo->prepare("SELECT * FROM users WHERE id = ?");
$stmt->execute([$id]);
$user = $stmt->fetch();

if (!$user) {
    header("Location: index.php");
    exit();
}

if (isset($_POST['submit']) && $_SERVER["REQUEST_METHOD"] == "POST")
{
    $formData = [];
    foreach ($_POST as $key => $value) 
    {
        $formData[$key] = trim($value);
    }
    
    $cin       = $formData["cin"] ?? '';
    $full_name = $formData["full_name"] ?? '';
    $tel       = $formData["tel"] ?? '';
    $email     = $formData["email"] ?? '';
    $birth     = $formData["birth"] ?? '';
    
    if (preg_match("/^[A-Z]{2}[0-9]{6}$/", $cin) == 0) 
    {
        exit("Invalid CIN format (e.g., AB123456)");
    }
    if (preg_match("/^[A-Za-z0-9.]+@[A-Za-z]+\.[a-z.]+$/", $email) == 0) 
    {
        exit("Invalid Email format");
    }
    if (preg_match("/^0[567]{1}[0-9]{8}$/", $tel) == 0) 
    {
        exit("Invalid Tel format (e.g., 0612345678)");
    }
    
    try 
    {
        $stmt = $pdo->prepare("UPDATE users SET cin = ?, full_name = ?, tel = ?, email = ?, birth = ? WHERE id = ?");
        $stmt->execute([$cin, $full_name, $tel, $email, $birth, $id]);
    
        header("Location: index.php?updated=1");
        exit();
    
    } catch (Throwable $e) {
        error_log("Unable to update user: " . $e->getMessage());
        exit("Unable to update the user. Please try again.");
    }
}
?>

<!DOCTYPE html>
<html lang="en">
<head>
    <meta charset="UTF-8">
    <meta name="viewport" content="width=device-width, initial-scale=1.0">
    <title>Modify User</title>
    <link rel="stylesheet" href="index.css"/>
</head>
<body>
    <h1>Modify User Infos</h1>
    <form method="POST">
        <label for="CIN">CIN :</label>
        <input type="text" name="cin" id="CIN" value="<?= htmlspecialchars($user['cin']) ?>" placeholder="AB123456" required>

        <label for="name">Full Name :</label>
        <input type="text" name="full_name" id="name" value="<?= htmlspecialchars($user['full_name']) ?>" required>

        <label for="tel">Tel :</label>
        <input type="text" name="tel" id="tel" value="<?= htmlspecialchars($user['tel']) ?>" placeholder="0612345678" required>

        <label for="email">Email :</label>
        <input type="email" name="email" id="email" value="<?= htmlspecialchars($user['email']) ?>" required>

        <label for="birth">Birthdate :</label>
        <input type="date" name="birth" id="birth" value="<?= htmlspecialchars($user['birth']) ?>" required>
        <br><br>
        <button class="update-user" type="submit" name="submit">Update User</button>
    </form>
</body>
</html>
\end{lstlisting}

\subsection{\texttt{delete.php}}
\begin{lstlisting}[language=PHP]
<?php
try {
    require_once "config.php";
} catch (Throwable $e) {
    error_log("Delete page database error: " . $e->getMessage());
    exit("The database is currently unavailable.");
}
if (isset($_GET["id"]))
{
    $id = $_GET["id"];

    try 
    {
        $stmt = $pdo->prepare("delete from users where id = ?");
        $stmt->execute([$id]);
        header("Location: index.php?deleted=1");
        exit();
    } 
    catch (Throwable $e) {
        error_log("Unable to delete user: " . $e->getMessage());
        exit("Unable to delete the user. Please try again.");
    }
}   
else 
{
    header("Location: index.php");
    exit();
}
?>
\end{lstlisting}

\subsection{\texttt{seed.php}}
\begin{lstlisting}[language=PHP]
<?php

require_once __DIR__ . "/vendor/autoload.php";
require_once __DIR__ . "/config.php";

$count = 15;
if (isset($argv[1])) {
    $count = filter_var($argv[1], FILTER_VALIDATE_INT, [
        "options" => ["min_range" => 1, "max_range" => 100],
    ]);
    if ($count === false) {
        fwrite(STDERR, "Usage: php seed.php [number of users, 1-100]\n");
        exit(1);
    }
}

$faker = Faker\Factory::create("fr_FR");
$insert = $pdo->prepare(
    "INSERT INTO users (cin, full_name, tel, email, birth) VALUES (?, ?, ?, ?, ?)"
);
$inserted = 0;

try {
    $pdo->beginTransaction();

    while ($inserted < $count) {
        $cin = strtoupper($faker->randomLetter()) . strtoupper($faker->randomLetter())
            . $faker->numerify("######");
        $birth = $faker->dateTimeBetween("-70 years", "-18 years")->format("Y-m-d");
        $tel = "0" . $faker->randomElement(["5", "6", "7"]) . $faker->numerify("########");

        try {
            $insert->execute([
                $cin,
                $faker->name(),
                $tel,
                $faker->unique()->safeEmail(),
                $birth,
            ]);
            $inserted++;
        } catch (PDOException $e) {
            if ($e->getCode() !== "23000") {
                throw $e;
            }
        }
    }

    $pdo->commit();
    echo "Inserted {$inserted} fake users.\n";
} catch (Throwable $e) {
    if ($pdo->inTransaction()) {
        $pdo->rollBack();
    }
    error_log("Unable to seed users: " . $e->getMessage());
    fwrite(STDERR, "Unable to seed users. Check the database configuration.\n");
    exit(1);
}
\end{lstlisting}

\subsection{\texttt{chart.php}}
\begin{lstlisting}[language=PHP]
<?php

try {
    require_once __DIR__ . "/config.php";
    require_once __DIR__ . "/vendor/autoload.php";
    \mitoteam\jpgraph\MtJpGraph::load("line");
    $stmt = $pdo->query("SELECT cin, birth FROM users ORDER BY id");
    $rows = $stmt->fetchAll();
} catch (Throwable $e) {
    error_log("Unable to load chart data: " . $e->getMessage());
    http_response_code(500);
    header("Content-Type: text/plain; charset=utf-8");
    exit("The chart data could not be loaded.");
}

if (!class_exists("Graph") || !class_exists("LinePlot")) {
    http_response_code(500);
    header("Content-Type: text/plain; charset=utf-8");
    exit("JpGraph is not installed. Run 'composer install' in the TP7 folder.");
}

$labels = [];
$birthYears = [];
foreach ($rows as $row) {
    $year = (int) date("Y", strtotime($row["birth"]));
    if ($year > 0) {
        $labels[] = $row["cin"];
        $birthYears[] = $year;
    }
}

if (empty($birthYears)) {
    $labels = ["No data"];
    $birthYears = [0];
}

$graph = new Graph(900, 420);
$graph->SetScale("textlin");
$graph->SetMargin(70, 30, 45, 100);
$graph->title->Set("CIN and birth year");
$graph->xaxis->title->Set("CIN");
$graph->yaxis->title->Set("Birth year");
$graph->xaxis->SetTickLabels($labels);
$graph->xaxis->SetLabelAngle(45);
$graph->SetColor("#ffffff");
$graph->SetBox(true);

$plot = new LinePlot($birthYears);
$plot->SetColor("#2563eb");
$plot->SetWeight(3);
$plot->mark->SetType(MARK_FILLEDCIRCLE);
$plot->mark->SetFillColor("#60a5fa");
$plot->mark->SetColor("#1d4ed8");
$graph->Add($plot);
$graph->Stroke();
\end{lstlisting}

\subsection{\texttt{index.css}}
\begin{lstlisting}[language=CSS]
* { box-sizing: border-box; }
body {
    font-family: Arial, sans-serif;
    max-width: 980px;
    margin: 35px auto;
    background: #111827;
    color: #f3f4f6;
    padding: 0 20px;
}
label 
{ 
    display: block; 
    margin: 12px 0 4px; 
} 
input, select, button 
{ 
    padding: 8px; 
    font-size: 1rem; 
}
.error { color: #fca5a5; }
.success { color: #86efac; }
.message { padding: 12px 16px; border-radius: 6px; margin-bottom: 20px; background: #1f2937; }
.muted { color: #9ca3af; }
.box 
{ 
    padding: 13px; background: #eff6ff; 
}
form {
    background: #1f2937;
    padding: 20px; 
    border: 1px solid #374151;
    border-radius: 6px;
    margin-bottom: 30px; 
}
table {
    width: 100%; 
    border-collapse: collapse; 
    background: #1f2937;
    border-radius: 6px;
    overflow: hidden;
    box-shadow: 0 1px 3px rgba(0,0,0,0.35);
}
th, td 
{ 
    padding: 12px; 
    border: 1px solid #374151;
    text-align: left; 
}
th { background: #374151; }
tr:nth-child(even) { background: #172033; }
input, select { background: #111827; color: #f3f4f6; border: 1px solid #4b5563; border-radius: 4px; }
.chart-card { margin-top: 30px; padding: 20px; background: #1f2937; border: 1px solid #374151; border-radius: 6px; }
.chart-card img { display: block; width: 100%; max-width: 900px; height: auto; margin-top: 15px; background: #fff; }
.action
{
    display: flex;
    gap: 5px;
}
.add-div
{
    display: flex;
    justify-content: space-between;
    padding-bottom: 20px;
    align-items: center;
}
.add-user
{
    height: 50px;
    background-color: #125dff;
    color: #eff6ff;
    border-radius: 10%;
}
.delete
{
    display: flex;
    justify-content: center;
    align-items: center;
    height: 50px;
    width: 75px;
    background-color: #ea3131;
    color: #eff6ff;
    border-radius: 10%;
    text-decoration: none;
}

.modify
{
    display: flex;
    justify-content: center;
    align-items: center;
    height: 50px;
    width: 75px;
    background-color: rgb(38, 187, 127);
    color: #eff6ff;
    border-radius: 10%;
    text-decoration: none;

}

.modify:hover, .delete:hover, .add-user-link:hover
{
    opacity: 0.8; 
}

.add-user-link
{
    display: flex;
    justify-content: center;
    align-items: center;
    height: 50px;
    width: 100px;
    background-color: #125dff;
    color: #eff6ff;
    border-radius: 10%;
    text-decoration: none;
}

.add-user-link:hover
{
    background-color: #3262d1; 
}
.update-user
{
    height: 50px;
    background-color: rgb(38, 187, 127);
    color: #eff6ff;
    border-radius: 10%;
}
\end{lstlisting}

\subsection{\texttt{composer.json}}
\begin{lstlisting}[language=JSON]
{
    "require": {
        "fakerphp/faker": "^1.24",
        "mitoteam/jpgraph": "^10.1"
    }
}
\end{lstlisting}
\end{document}
